% Part:sets-functions-relations
% Chapter: size-of-sets
% Section: comparing-sizes

\documentclass[../../../include/open-logic-section]{subfiles}

\begin{document}

\olfileid{sfr}{siz}{car}

\olsection{જુદાં કદના ગણો અને કૅન્ટૉરનું પ્રમેય}

\begin{explain}
બે ગણોનું કદ સરખું હોવાના વિચારને આપણે ચોક્કસ રીતે રજૂ કર્યો છે.
એક ગણ બીજા કરતાં નાનો હોવાના વિચારને પણ ચોક્કસ રીતે રજૂ કરી શકાય.
“નાનો (અથવા સામ્ય ગણ)” હોવાની આપણી વ્યાખ્યામાં બંને ગણો વચ્ચેના
!!a{bijection}ને બદલે પહેલા ગણમાંથી બીજા ગણમાં જતા !!a{injection}ની
જરૂર પડશે. આવું વિધેય અસ્તિત્વમાં હોય, તો પહેલા ગણનું કદ બીજા
ગણના કદ કરતાં ઓછું અથવા બરાબર છે. સહજ રીતે, એક ગણમાંથી બીજા ગણમાં
જતો !!a{injection} ખાતરી આપે છે કે વિધેયના વિસ્તારમાં તેના પ્રદેશ
જેટલા તો !!{element}s છે જ, કારણ કે પ્રદેશના કોઈ બે !!{element}sનું
નિરૂપણ વિસ્તારના એક જ !!{element} પર થતું નથી.
\end{explain}

\begin{defn}
$A$ એ $B$ કરતાં \emph{મોટો નથી}, જે $\cardle{A}{B}$ લખીએ,
જો અને માત્ર જો કોઈ !!a{injection} $f \colon A \to B$ હોય.
\end{defn}

સ્પષ્ટ છે કે આ સંબંધ સ્વવાચક અને પરંપરિત છે, પરંતુ સંમિત નથી
(આ સાબિત કરવાનું અભ્યાસ તરીકે છોડ્યું છે). એક ગણ બીજા કરતાં
(ચુસ્તપણે) નાનો હોવાનું જણાવતો ખ્યાલ પણ રજૂ કરી શકાય છે.

\begin{defn}
$A$ એ $B$ કરતાં \emph{નાનો} છે, જે $\cardless{A}{B}$ લખીએ,
જો અને માત્ર જો કોઈ !!a{injection}~$f\colon A \to B$ હોય પરંતુ
કોઈ !!{bijection}~$g\colon A \to B$ ન હોય; એટલે કે,
$\cardle{A}{B}$ અને $\cardneq{A}{B}$.
\end{defn}

સ્પષ્ટ છે કે આ સંબંધ અસ્વવાચક અને પરંપરિત છે. (આ સાબિત કરવાનું
અભ્યાસ તરીકે છોડ્યું છે.) આ સંજ્ઞા વાપરીને કહી શકાય કે ગણ $A$ એ
!!{enumerable} હોય જો અને માત્ર જો $\cardle{A}{\Nat}$, અને $A$ એ
!!{nonenumerable} હોય જો અને માત્ર જો $\cardless{\Nat}{A}$.
આથી
\oliflabeldef{sfr:siz:nen-alt:thm:nonenum-pownat}{%
\olref[sfr][siz][nen-alt]{thm:nonenum-pownat}ને
$\cardless{\Nat}{\Pow{\Nat}}$ એવા અવલોકન તરીકે}{%
\olref[sfr][siz][nen]{thm:nonenum-pownat}ને
$\cardless{\PosInt}{\Pow{\PosInt}}$ એવા અવલોકન તરીકે} ફરી રજૂ
કરી શકાય છે. હકીકતમાં, \citet{Cantor1892}એ સાબિત કર્યું કે આ છેલ્લી
વાત \emph{સર્વથા સામાન્ય} છે:

\begin{thm}[કૅન્ટૉર]\ollabel{thm:cantor}
કોઈ પણ ગણ $A$ માટે, $\cardless{A}{\Pow{A}}$.
\end{thm}

\begin{proof}
નિરૂપણ $f(x) = \{x\}$ એ !!a{injection} $f \colon A \to \Pow{A}$
છે, કારણ કે $x \neq y$ હોય, તો ઘટકો દ્વારા નિર્ધારિત સમાનતાના
સિદ્ધાંતથી $\{x\} \neq \{y\}$ પણ હોય, અને તેથી $f(x) \neq f(y)$.
આથી $\cardle{A}{\Pow{A}}$.

\begin{editorial}
\olref[nen]{sec} સમાવેલો હોય તો આપણે વિગતવાર સાબિતી રજૂ કરીએ છીએ;
અન્યથા \olref[nen-alt]{sec}ને અનુરૂપ ટૂંકી સાબિતી રજૂ કરીએ છીએ.
\end{editorial}

\oliflabeldef{sfr:siz:nen:sec}{%
  હવે આપણે બતાવીશું કે !!a{surjective} વિધેય~$g\colon A \to
  \Pow{A}$ અસ્તિત્વમાં હોઈ જ ન શકે, તેથી !!a{bijective} વિધેય તો
  અસ્તિત્વમાં ન જ હોય, અને આમ $\cardneq{A}{\Pow{A}}$. વિરોધાર્થે,
  ધારો કે $g\colon A \to \Pow{A}$. $g$ પૂર્ણ હોવાથી દરેક
  $x \in A$નું નિરૂપણ કોઈ ઉપગણ $g(x) \subseteq A$ પર થાય છે.
  આપણે બતાવી શકીએ કે $g$ વ્યાપ્ત ન હોઈ શકે. તે માટે આપણે એવો
  ઉપગણ~$\overline{A} \subseteq A$ વ્યાખ્યાયિત કરીએ છીએ જે
  વ્યાખ્યા પ્રમાણે $g$ના વિસ્તારમાં હોઈ ન શકે. ધારો કે
  \[
  \overline{A} = \Setabs{x \in A}{x \notin g(x)}.
  \]
  દરેક $x \in A$ માટે $g(x)$ વ્યાખ્યાયિત હોવાથી $\overline{A}$
  સ્પષ્ટ રીતે $A$નો સુવ્યાખ્યાયિત ઉપગણ છે. પરંતુ તે $g$ના વિસ્તારમાં
  હોઈ ન શકે. મનસ્વી $x \in A$ લો; આપણે બતાવીશું કે
  $\overline{A} \neq g(x)$. જો $x \in g(x)$, તો તે
  $x \notin g(x)$ને સંતોષતો નથી, અને તેથી $\overline{A}$ની વ્યાખ્યા
  પ્રમાણે $x \notin \overline{A}$. જો $x \in \overline{A}$, તો
  તેણે $\overline{A}$નો વ્યાખ્યાયી ગુણધર્મ સંતોષવો જ પડે; એટલે કે,
  $x \in A$ અને $x \notin g(x)$. $x$ મનસ્વી હોવાથી, આ બતાવે છે કે
  દરેક $x \in A$ માટે $x \in g(x)$ જો અને માત્ર જો
  $x \notin \overline{A}$, અને તેથી $g(x) \neq \overline{A}$.
  બીજા શબ્દોમાં, $\overline{A}$ એ $g$ના વિસ્તારમાં હોઈ ન શકે, જે
  $g$ વ્યાપ્ત હોવાની ધારણાનો વિરોધાભાસ છે.
  \sourcecorrection{OLSIZ-007}{મૂળની \texttt{comparing-size.tex},
  પંક્તિઓ 91--95માં મનસ્વી \(x\in A\) માટેના નિષ્કર્ષને ભૂલથી
  “દરેક \(x\in\overline A\) માટે” મર્યાદિત કર્યો છે. અહીં દલીલને
  નિયંત્રિત કરતી મનસ્વી પસંદગી મુજબ \(x\in A\) લખ્યું છે.}}{બાકી
  $\cardneq{A}{\Pow{A}}$ બતાવવાનું છે. વિરોધ દ્વારા સાબિતી માટે
  ધારો કે $\cardeq{A}{\Pow{A}}$; એટલે કે કોઈ
  !!{bijection} $g \colon A \to \Pow{A}$ છે. હવે આ ગણ જુઓ:
  \[
    D = \Setabs{x \in A}{x \notin g(x)}
  \]
  નોંધો કે $D \subseteq A$, તેથી $D \in \Pow{A}$. $g$ એ
  !!a{bijection} હોવાથી કોઈ $y \in A$ માટે $g(y) = D$. પરંતુ હવે:
  \[
    y \in g(y) \text{ જો અને માત્ર જો } y \in D \text{ જો અને માત્ર જો } y \notin g(y).
  \]
  આ વિરોધાભાસ છે; તેથી $\cardneq{A}{\Pow{A}}$.}{}
\end{proof}

\begin{explain}
\oliflabeldef{sfr:siz:nen:thm:nonenum-pownat}{\olref{thm:cantor}ની
  સાબિતીને \olref[nen]{thm:nonenum-pownat}ની સાબિતી સાથે સરખાવવી
  બોધપ્રદ છે. ત્યાં આપણે બતાવ્યું હતું કે $\PosInt$ના ઉપગણોની કોઈ પણ
  યાદી $Z_1$, $Z_2$, \dots, માટે એવો સંખ્યાઓનો ગણ~$\overline{Z}$
  બનાવી શકાય જે યાદીમાં ન હોવાની ખાતરી હોય. તે યાદીમાં નથી તેની
  ખાતરીનું કારણ એ હતું કે દરેક $n \in \PosInt$ માટે $n \in Z_n$
  જો અને માત્ર જો $n \notin \overline{Z}$. આ રીતે $Z_n$ અને
  $\overline{Z}$ પૈકી એકનો !!a{element} હોય પરંતુ બીજાનો ન હોય એવી
  કોઈ સંખ્યા હંમેશાં મળે છે. અહીં આપણે એ જ વિચારને અનુસરીએ છીએ,
  પરંતુ હવે સૂચકો~$n$ એ $\PosInt$ના બદલે $A$ના !!{element}s છે.
  ગણ $\overline{A}$ને એવી રીતે વ્યાખ્યાયિત કર્યો છે કે તે દરેક
  $x \in A$ માટે $g(x)$થી જુદો હોય, કારણ કે $x \in g(x)$ જો અને
  માત્ર જો $x \notin \overline{A}$. ફરી, $A$નો એવો !!a{element}
  હંમેશાં હોય છે જે $g(x)$ અને $\overline{A}$ પૈકી એકનો
  !!a{element} હોય પણ બીજાનો ન હોય. અને જેમ $\overline{Z}$ યાદી
  $Z_1$, $Z_2$, \dots, માં હોઈ શકતો નથી, તેમ $\overline{A}$ એ
  $g$ના વિસ્તારમાં હોઈ શકતો નથી.}{}

\oliflabeldef{sfr:siz:nen-alt:thm:nonenum-pownat}{\olref{thm:cantor}ની
  સાબિતીને \olref[nen-alt]{thm:nonenum-pownat}ની સાબિતી સાથે સરખાવવી
  બોધપ્રદ છે. ત્યાં આપણે બતાવ્યું હતું કે $\Nat$ના ઉપગણોની કોઈ પણ
  યાદી $N_0$, $N_1$, $N_2$, \dots, માટે એવો સંખ્યાઓનો ગણ~$D$
  બનાવી શકાય જે યાદીમાં ન હોવાની ખાતરી હોય. તેનું કારણ એ હતું કે
  દરેક $n \in \Nat$ માટે $n \in N_n$ જો અને માત્ર જો $n \notin D$.
  અહીં આપણે એ જ વિચારને અનુસરીએ છીએ, પરંતુ હવે સૂચકો~$n$ એ $\Nat$ના
  બદલે $A$ના !!{element}s છે. ગણ $D$ને એવી રીતે વ્યાખ્યાયિત કર્યો છે
  કે તે દરેક $x \in A$ માટે $g(x)$થી જુદો હોય, કારણ કે
  $x \in g(x)$ જો અને માત્ર જો $x \notin D$.}{}

આ સાબિતીની રસેલના વિરોધાભાસની સાબિતી,
\olref[sfr][set][rus]{thm:russells-paradox}, સાથેની તુલના પણ કરવા
જેવી છે. ખરેખર, કૅન્ટૉરનું પ્રમેય રસેલના પોતાના વિરોધાભાસની પ્રેરણા હતું.
\end{explain}

\begin{prob}
  બતાવો કે કોઈ પણ ગણ~$A$ માટે !!a{injection}
  $g\colon \Pow{A} \to A$ હોઈ ન શકે. સંકેત: ધારો કે
  $g\colon \Pow{A} \to A$ એ !!{injective} છે. ગણ
  $D = \Setabs{g(B)}{B \subseteq A \text{ અને } g(B) \notin B}$
  વિચારો. $x = g(D)$ લો. વિરોધાભાસ મેળવવા માટે $g$ !!{injective}
  હોવાની હકીકત વાપરો.
\end{prob}

\end{document}
